% Lösungen Einheit 4 von 5 – Dezimalzahlen multiplizieren und dividieren (zusammenbau v0.1)
\einheitenkopf{Einheit 4 von 5 · Dezimalzahlen multiplizieren und dividieren}

\erg{23}{a) 2 Kommastellen}
%% TODO Lösung der Erklärzeile 24f) fehlt in der Bank
\erg{24}{a) $37{,}2$ \quad b) $58$ \quad c) $450$ \quad d) $0{,}93$ \quad e) $126$ \quad f) \ldots \quad g) $4{,}57$ \quad h) $6{,}9$ \quad i) $0{,}84$ \quad j) $0{,}08$ \quad k) $2{,}1$ \quad l) $56 : 7 = 8$ \quad m) $25\,000 : 100 \cdot 8 = 2\,000$ l; $2\,000 \cdot 1{,}624\text{ €} = 3\,248\text{ €}$ \quad n) $0{,}2^2 = 0{,}04$}
\erg{25}{a) Sie hat die Kommastellen nicht gezählt; beide Faktoren haben je eine, das Ergebnis braucht zwei. Richtig: $0{,}7 \cdot 0{,}4 = 0{,}28$. \quad b) Jede Ziffer wird zehnmal so viel wert: aus Zehnteln werden Einer, aus Hundertsteln Zehntel. \quad c) $1{,}5 \cdot 2{,}40\text{ €} = 3{,}60\text{ €}$; der Beutel ist günstiger.}
